% A PREFIX READS WHAT IT PREFIXES ITSELF, AND SO HIDES IT FROM THE TRACE.
%
% \tracingcommands draws a line for the prefix and nothing for the
% command it prefixes:
%
%   \global\count5=1          {\global}
%   \global\advance\count5by1 {\global}
%   \long\def\aa{}            {\long}
%   \outer\def\bb{}           {\outer}
%   \global\let\cc=\relax     {\global}
%   \global\long\def\dd{}     {\global}      -- not the \long either
%   \global\relax\count5=2    {\global}      -- nor a \relax between
%   \global \count5=3         {\global}      -- nor the space
%   \global\setbox1=\hbox{}   {\global}
%   \global\futurelet\ee\relax{\global}
%
% The same shape as \immediate, which hides \write, \openout and
% \closeout: a command read by something other than the main loop is not
% a command the main loop obeyed. See
% tests/trip/probes/what-mode-a-write-text-is-expanded-in.tex.
%
% A PREFIX WITH SOMETHING IT CANNOT PREFIX draws two lines, because the
% command is put back and met again once the prefix has been given up --
% the same way \immediate\relax does.
\tracingonline=1 \tracingcommands=2
\message{[a]}\global\count5=1
\message{[b]}\global\advance\count5 by1
\message{[c]}\long\def\aa{}
\message{[d]}\outer\def\bb{}
\message{[e]}\global\let\cc=\relax
\message{[f]}\global\relax\count5=2
\message{[g]}\global\long\def\dd{}
\message{[h]}\global\setbox1=\hbox{}
\message{[i]}\global\futurelet\ee\relax
\message{[j]}
\message{[k]}\global \count5=3
\message{[l]}
\tracingcommands=0
\end
